\documentclass[headsepline,11pt]{scrbook}
\usepackage[a4paper]{geometry}
\usepackage{amsmath,amsfonts,amsthm,amssymb}
\usepackage[ngerman]{babel}
%
\newtheorem{lem}{Lemma}[section]
\newtheorem{thm}[lem]{Satz}
\theoremstyle{definition}
\newtheorem{df}[lem]{Definition}
%
\def\R{\mathbb R}
\def\C{\mathbb C}
\DeclareMathOperator{\vol}{vol}
%
\usepackage{tikz}
%
\begin{document}
%
%%%%% KEiNE AENDERUNGEN UEBER DIESER ZEILE %%%%%%%%%%%%%%%%%%%%%%%%%%%%%%%%%%%%%
%%%%%%%%%%%%%%%%%%%%%%%%%%%%%%%%%%%%%%%%%%%%%%%%%%%%%%%%%%%%%%%%%%%
%
\setcounter{chapter}{7}
\chapter{W\"armeleitungskerne }

Ausgearbeitet von: {\sffamily\bfseries Philipp Beck}\\


\section{Einleitung}
Unser Ziel ist es, die W\"armleitungsgleichung 
\begin{equation}\label{eq:heat_equation}
\begin{split}
\frac{\partial u}{\partial  \mathrm{t}} &= - \mathcal{L} u\\
u(0,\cdot) &= u_0
\end{split}
\end{equation}
auf Graphen und Teilgraphen zu l\"osen.
%Dies werden wir auf Graphen und Teilgraphen unter Randbedingungen mithilfe des W\"armeleitungskerns bewerkstelligen.
Die L\"osung werden wir durch
\begin{align}
u(t, \cdot) = H_t u_0
\end{align}
erhalten. Hierbei ist $ H_t $ der W\"armeleitungsoperator.
Der Laplaceoperator ist auf einen Graphen bzw. induzierten Teilgraphen unter Randbedingungen selbstadjungiert.
Nach dem Spektralsatz von Riesz-Schauder erhalten wir eine Orthonormalbasis aus Eigenfunktionen und die Darstellung
\begin{equation}\label{eq:representation_laplace}
\mathcal{L}
= 
\sum \limits_{\lambda \in \sigma(\mathcal{L})}
\lambda  \ I_\lambda
\end{equation}
\"uber die Eigenwerte $ \lambda \in \sigma(\mathcal{L}) $ und die zugeh\"origen Eigenprojektionen $ I_\lambda $.
Desweiteren seien die Eigenfunktionen im allgemeinen  
reell gew\"ahlt.

\begin{lem}\label{th:representation_laplace}
	F\"ur $ k \in \mathbb{N} $ gilt
	$ \mathcal{L}^k = \sum_\lambda  \lambda^k  I_\lambda $.
\end{lem}

\begin{proof}
	Der Induktionsanfang ist klar. 
	Angenommen die Aussage gilt f\"ur ein $ k $.
	Dann folgt mit \eqref{eq:representation_laplace} und $ I_{\lambda^\prime} I_\lambda  f  = 0 $ 
	%f\"ur $ \lambda \neq \lambda^\prime $
	\begin{align*}
	\mathcal{L}^{k+1} f
	= 
	\left( \sum \limits_{\lambda \in \sigma(\mathcal{L})} \lambda^k I_\lambda \right)
	\left( \sum \limits_{\lambda^\prime \in \sigma(\mathcal{L})} \lambda^\prime  I_\lambda^\prime \right) f
	=
	\sum \limits_{\lambda} \lambda^{k+1} I_\lambda f +
	\sum \limits_{\lambda \neq \lambda^\prime} 
	\lambda^k \lambda^\prime 
	\underbrace{I_\lambda I_{\lambda^\prime } f}_{=0}
	=
	\left(\sum \limits_{\lambda} \lambda^{k+1} I_\lambda \right) f
	\end{align*}
	f\"ur $ f \in L^2(G) $.
\end{proof}



\section{Der W\"armeleitungskern auf dem Gesamtgraphen }
Sei $ G  $ ein Graph mit endlich vielen Knoten.
%(\textcolor{red}{\textbf{oder induzierte Subgraphen etc.}})
Wie schon erw\"ahnt, erhalten wir mit dem Spektralsatz von Riesz-Schauder
eine ONB aus reellen Eigenfunktionen f\"ur $ L^2(G) $.
Durch
\begin{equation}
\begin{split}\label{eq:kernel_projection}
I_\lambda f(x) &=
\sum \limits_{\phi \in E_\lambda} \langle f , \phi \rangle \phi(x)
=
\sum \limits_{\phi \in E_\lambda} 
\left(\sum \limits_{y \in V(G)} 
f(y) \phi(y)  \mathrm{d}_y\right) \phi(x)\\
&=
\sum \limits_{y \in V(G)} f(y) \mathrm{d}_y  
\underbrace{\left(\sum \limits_{\phi \in E_\lambda} \phi(x) \phi(y)\right) }_{i_\lambda(x,y)}
=
\langle f, i_\lambda(x,\cdot) \rangle_G
\end{split}
\end{equation}
erhalten wir eine Darstellung f\"ur die Projektion auf den Eigenraum $ E_\lambda $. 
Wir bezeichnen $ i_\lambda $ als \glqq Integralkern\grqq~ von $ I_\lambda $ und verwenden diesen, um den W\"armeleitungskern herzuleiten.

\begin{lem}\label{th:projection_kernel_different_eigenvalues}
	F\"ur $ \lambda, \lambda^\prime \in \sigma(\mathcal{L}) $ mit $ \lambda \neq \lambda^\prime $ gilt
	\begin{align}\label{eq:projection_kernel_different_eigenvalues}
%	\sum \limits_{y \in V(G)} i_\lambda(x,y) i_{\lambda^\prime}(y,z) \mathrm{d}_y = 0
	\langle i_\lambda(x,\cdot), i_{\lambda^\prime}(\cdot,z) \rangle_G
	=0
	\end{align}
	f\"ur beliebige  $ x,z \in V(G) $.
\end{lem}

\begin{proof}
	Es gilt:
	\begin{align*}
	\langle i_\lambda(x,\cdot), i_{\lambda^\prime}(\cdot,z) \rangle_G
	&=
	\sum \limits_{y \in V(G)}
	\left( \sum \limits_{\phi_1 \in E_\lambda} \phi_1(x) \phi_1(y) \right)
	\left(\sum \limits_{\phi_2 \in E_{\lambda^\prime}} \phi_2(y) \phi_2(z)\right) \mathrm{d}_y\\
	&=
	\sum \limits_{y \in V(G) }
	\sum \limits_{\substack{\phi_1 \in E_\lambda\\
					\phi_2 \in E_{\lambda^\prime}
			}}
			\phi_1(x) \phi_1(y) \phi_2(y) \phi_2(z)
			\mathrm{d}_y\\
	&=
	\sum \limits_{\substack{\phi_1 \in E_\lambda\\
			\phi_2 \in E_{\lambda^\prime}
			}}
		\phi_1(x) \phi_2(z) \langle \phi_1,\phi_2 \rangle
	= 0.	
	\end{align*}
\end{proof}

\begin{lem}\label{th:projection_kernel_same_eigenvalues}
	F\"ur $ \lambda \in \sigma(\mathcal{L}) $ gilt
	\begin{align}\label{eq:projection_kernel_same_eigenvalues}
%	\sum \limits_{y \in V(G)} 
%	i_\lambda(x,y) i_\lambda(y,z) \mathrm{d}_y
	\langle i_\lambda(x,\cdot) , i_\lambda(\cdot,z) \rangle_G
	=
	i_\lambda(x,z)
	\end{align}
	f\"ur $ x,z \in V(G) $.
	
\end{lem}

\begin{proof}
	Zun\"achst gilt: 
	\begin{align*}
%	\sum \limits_{y \in V(G)} 
%	i_\lambda(x,y) i_\lambda(y,z) \mathrm{d}_y
	\langle i_\lambda(x, \cdot), i_\lambda(\cdot, z) \rangle_G
	=
	\sum \limits_{y \in V(G)} 
	\left( \sum \limits_{\phi \in E_\lambda} \phi(x) \phi(y) \right)
	\left(\sum \limits_{\psi \in E_\lambda} \psi(y) \psi(z)\right) \mathrm{d}_y.
	\end{align*}
	Mit $ \langle \phi , \psi \rangle = 0  $ f\"ur $ \phi \neq \psi $ erhalten wir
	\begin{align*}
	\langle i_\lambda(x, \cdot), i_\lambda(\cdot, z) \rangle_G
	=
	\sum \limits_{\phi \in E_\lambda} 
	\phi(x)\phi(z) \| \phi\|
	=
	i_\lambda(x,z)
	\end{align*}
\end{proof}

\begin{df}
Sei $ t \in [0,\infty) $.
Der W\"armeleitungsoperator 
$ H_t $ ist definiert durch
\begin{align}\label{eq:definition_heatkernel}
H_t : L^2(G) \to L^2(G),
u_0 \mapsto u(t,\cdot).
\end{align}
Hierbei ist $ u $ die L\"osung der W\"armeleitungsgleichung 
\eqref{eq:heat_equation}.
\end{df}
\begin{thm}\label{th:repres_heat_kernel_global}
	F\"ur den W\"armeleitungsoperator gilt
	\begin{align}\label{eq:heat_kernel_def_whole}
	u(t,x) = (H_t u_0)(x) = 
	\sum \limits_{\lambda \in \sigma(\mathcal{L})} e^{- \lambda t } (I_\lambda u_0) (x) 
	=
	\langle u_0,  h_t(x,\cdot) \rangle_G.
	%\sum \limits_{y \in V(G)} h_t(x,y) f(y) \mathrm{d}_y.
	\end{align}
	f\"ur die Anfangsbedingung $ u_0 \in L^2(G) $.
	%Insbesondere gilt $ u(0,x) = u_0(x) $.
	Hierbei ist
	\begin{align}
	h_t(x,y) := \sum \limits_{\lambda \in \sigma(\mathcal{L})}
	e^{- \lambda t} i_\lambda(x,y)
	\end{align}
	der W\"armeleitungskern auf $ G $.
\end{thm}

\begin{proof}
	Wir betrachten
	\begin{align*}
	e^{-t \mathcal{L}}
	=
	\sum \limits_{k=0}^\infty
	\frac{(-t)^k}{k!} \mathcal{L}^k
	\overset{\ref{th:representation_laplace}}{=}
	\sum \limits_{\lambda \in \sigma(\mathcal{L})}
	\sum \limits_{k = 0}^\infty \frac{(-t \lambda)^k}{k!} I_\lambda
	=
	\sum \limits_{\lambda \in \sigma(\mathcal{L})} e^{-t \lambda } I_\lambda.
	\end{align*}
	Wegen 
	\begin{align*}
	\partial_t e^{-t \mathcal{L}} = - \mathcal{L} e^{-t \mathcal{L}}
	\end{align*}
	erf\"ullt $ H_t u_0 := e^{-t \mathcal{L}} u_0 $ die W\"armeleitungsgleichung \eqref{eq:heat_equation}.
	Damit folgt durch die Gleichung \eqref{eq:kernel_projection} die Darstellung:
	\begin{align*}
	(H_t u_0)(x) &= 
	\sum \limits_{\lambda \in \sigma(\mathcal{L})} e^{- \lambda t } (I_\lambda u_0) (x) 
	=
	\sum \limits_{\lambda \in \sigma(\mathcal{L})} e^{- \lambda t }
	\left(\sum \limits_{y \in V(G)} u_0(y) \mathrm{d}_y i_\lambda(x,y) \right)\\
	&=
	\sum \limits_{y \in V(G)} u_0(y) \left(\sum \limits_{\lambda \in \sigma(\mathcal{L})} e^{- \lambda t } i_\lambda(x,y) \right) \mathrm{d}_y\\
	&=
	\sum \limits_{y \in V(G)} u_0(y) h_t(x,y) \mathrm{d}_y
	=
	\langle h_t(x,\cdot) ,u_0 \rangle_G.
	\end{align*}
\end{proof}
Wegen $ 0 \in \sigma(\mathcal{L}) $ und $ \mathbf{1} \in E_0 $ gilt aufgrund der Gleichung \eqref{eq:heat_kernel_def_whole}
die Identit\"at\\
$ (H_t \textbf{1})(x) = \textbf{1}(x) $.

\begin{lem}
	F\"ur $ 0 \leq s \leq t $ gilt:
	\begin{align}
%	\sum 
%	\limits_{y \in V(G)}
%	h_s(x,y) h_{t-s}(y,z) 
%	\mathrm{d}_y
	\langle h_s(x,\cdot), h_{t-s}(\cdot,z) \rangle_G
	=
	h_t(x,z).
	\end{align}
\end{lem}

\begin{proof}
	Durch Einsetzen erhalten wir:
	\begin{align*}
	&\sum 
	\limits_{y \in V(G)}
	h_s(x,y) h_{t-s}(y,z) 
	\mathrm{d}_y
	=
	\sum \limits_{y \in V(G)}
	\left(\sum \limits_{\lambda \in \sigma(\mathcal{L})}
	e^{-\lambda s}
	i_\lambda(x,y)
	\right)
	\left(\sum \limits_{\lambda^\prime \in \sigma(\mathcal{L})}
	e^{-\lambda^\prime (t-s)}
	i_\lambda(y,z)
	\right)
	\mathrm{d}_y\\
	=
	&\sum \limits_{y \in V(G)}
	\left(
	\sum \limits_{\lambda \in \sigma(\mathcal{L})}
	e^{t\lambda}
	i_\lambda(x,y) i_\lambda(y,z)
	+
	\sum \limits_{\lambda \neq \lambda^\prime}
	e^{-\lambda s} e^{- \lambda^\prime(t-s)}
	i_\lambda(x,y)i_{\lambda^\prime}(y,z)
	\right)
	\mathrm{d}_y\\
	\overset{\ref{th:projection_kernel_different_eigenvalues}}{=}
	&\sum \limits_{\lambda\in \sigma(\mathcal{L})}
	e^{t \lambda}
	\left(\sum \limits_{y \in V(G)} 
	i_\lambda(x,y) i_\lambda(y,z) \mathrm{d}_y
	\right)
	\overset{\ref{th:projection_kernel_same_eigenvalues}}{=}
	\sum \limits_{\lambda \in \sigma(\mathcal{L})}
	e^{-\lambda t} i_\lambda(x,z)
	= h_t(x,z)
	\end{align*}
\end{proof}

\begin{thm}
	F\"ur $ 0 \leq s \leq t $ gilt
	\begin{align}
	H_{s} H_{t-s} = H_t.
	\end{align}
\end{thm}

\begin{proof}
	Der Beweis folgt direkt durch Anwenden des letzten Lemmas.
	Es gilt:
	\begin{align*}
	(H_s H_{t-s} f)(x)
	&=
	\sum_y h_s(x,y) (H_{t-s} f)(y) \mathrm{d}_y
	=
	\sum_z \left(\sum_y h_s(x,y) h_{t-s}(y,z) \mathrm{d}_y \right) f(z) \mathrm{d}_z\\ 
	&= (H_t f)(x).
	\end{align*}
\end{proof}


\begin{thm}\label{th:pos_heat_kernel_global}
	Es gilt 
	\begin{align}
	h_t(x,y) \geq 0
	\end{align}
	f\"ur $ x,y \in G $.
\end{thm}

\begin{proof}
	Sei $ u_0 : V(G)  \to [0,\infty]$ beliebig.
	Dann existiert f\"ur ein beliebiges $ t \in [0, \infty) $ ein 
	$ x_{\textrm{max}}  \in  G$ mit
	\begin{align*}
	(H_t u_0)(x_{\mathrm{max}}) = u(t,x_{\mathrm{max}}) = \max\limits_{x \in V(G) } u(t,x).
	\end{align*}
	F\"ur den Laplace erhalten wir dann
	\begin{align*}
	\mathcal{L}u(t,x_{\mathrm{max}} ) =
	\frac{1}{\mathrm{d}_{x_{\mathrm{max}}}}
	\sum 
	\limits_{y \sim x_{\mathrm{max}}}
	(u(t,x_\mathrm{max}) - u(t,y)) \geq 0
	\end{align*}
	Wegen der W\"armeleitungsgleich \eqref{eq:heat_equation} folgt dann auch $ \partial_\mathrm{t} u(t,x_{\mathrm{max}}) \leq 0 $.
	Durch ein analoges Argument f\"ur das Minimum gilt $ \partial_\mathrm{t} u(t,x_\mathrm{min}) \geq 0 $.
	Damit wird insbesondere das Maximum mit der Zeit h\"ochstens kleiner und das Minimum gr\"o\ss{}er. 
	Damit gilt wegen $ u_0 \geq 0  $ auch $ (H_t u_0) \geq 0 $.
	Mit 
	\begin{align*}
	u_z(y) =
	\begin{cases}
	\frac{1}{\mathrm{d}_z} , &\quad z=y\\
	\ 0 ,						&\quad z \neq y
	\end{cases} 
	\end{align*}
	erhalten wir
	\begin{align*}
	(H_t u_z )(x) = \sum \limits_{y \in V(G)}
	u_z(y) h_t(x,y) \mathrm{d}_y
	= h_t(x,z) \geq 0
	\end{align*}
	f\"ur alle $ x,z \in G $.
	Aufgrund der Symmetrie von $ h_t $ folgt die gew\"unschte Aussage.
\end{proof}

Als Folgerung ergibt sich f\"ur $ u_0 \geq v_0 $ direkt
\begin{align}
(H_t u_0)(x)
=
\sum \limits_{y \in V(G)}
u_0(y) h_t(x,y) \mathrm{d}_y
\geq 
\sum \limits_{y \in V(G)}
v_0(y) h_t(x,y) \mathrm{d}_y
=
(H_t v_0)(x)
\end{align}
f\"ur alle $ x \in V(G) $.
Fassen wir den W\"armeleitungsoperator $ H_t  $ als Operator von $ L^\infty $ nach $ L^\infty $ auf, so gilt f\"ur die Operatornorm $ \| H_t f \|_\infty = 1 $.
Dies folgt wegen
\begin{align*}
\| H_t f \|_\infty
= 
\max \limits_{ \| f \|_\infty = 1}
\| H_t f \|_\infty
\leq 
\max \limits_{ \| f \|_\infty = 1} 
\max \limits_{x \in V(G)}
| (H_t |f |)(x) | 
\leq 
\max \limits_{x \in V(G)} 
| \underbrace{(H_t \textbf{1})}_{
	= \textbf{1}(x)
	}(x) |
= 1.
\end{align*}

%\begin{thm}
%	Sei $ x \in V(G) $.
%	Dann gilt
%	\begin{align}
%	\sum \limits_{y \in V(G)}
%	h_t(x,y) d_y  = \textbf{1}(x).
%	\end{align}
%\end{thm}
%
%\begin{proof}
%%	Wir betrachten die Abbildung
%%	\begin{align*}
%%	f : G \to \mathbb{R}, \
%%	x \mapsto \mathrm{d}_x.
%%	\end{align*}
%	Es gilt
%	\begin{align*}
%	(H_t \textbf{1})(x) = \sum \limits_y h_t(x,y) \mathrm{d}_y
%	\end{align*}
%	und wegen $ 0 \in \sigma(\mathcal{L}) $ folgt
%	\begin{align*}
%	(H_t f)(x)
%	=
%	\sum \limits_{\lambda \in \sigma(\mathcal{L})}
%	e^{-\lambda t} I_\lambda f(x)
%	=
%	\mathrm{d}_x e^0 \textbf{1}(x) = \mathrm{d}_x
%	\end{align*}
%	und damit die Aussage.
%\end{proof}

\begin{thm}\label{th:L_2_estimate_global}
	Sei $ f \in L^2(G) $.
	Dann gilt
	\begin{align*}
	\| H_t f - \overline{f}  \|_2 \leq e^{-\lambda_1 t} \| f \|_2. 
	\end{align*}
	Hierbei bezeichnet 
	\begin{align}
	\overline{f} := 
	\frac{1}{\vol(G)} \langle f ,\textbf{1} \rangle
	\end{align}
	den Mittelwert \"uber $ G $.
\end{thm}
\begin{proof}
	Nach Definition gilt $ \overline{f} = I_0 f $.
	Aus dem Satz \ref{th:repres_heat_kernel_global} folgt
	\begin{align*}
	(H_t f)(x) - (I_0 f)(x) = 
	\sum \limits_{\lambda \in \sigma(\mathcal{L}) \setminus \{0\}}
	e^{-\lambda t} (I_\lambda f)(x)
	=
	\sum \limits_{i = 1}^{n-1} e^{-\lambda_i t} (I_{\lambda_i} f)(x)
	\end{align*}
	mit $ 0 < \lambda_1 < ... < \lambda_{n-1} $.
	Damit erhalten wir durch
	\begin{align*}
	\| H_t f - I_0 f  \|_2^2 
	&=
	\sum \limits_{x \in V(S)} \left| 
	\sum \limits_{i=1}^{n-1} e^{-\lambda_i t} I_{\lambda_i} f(x)
	\ \right|^2 \mathrm{d}_x
	\leq
	\sum \limits_{i=1}^n e^{-\lambda_i 2 t}  
	\left(\sum \limits_x | I_{\lambda_i} f(x) |^2 \mathrm{d}_x\right)\\
	&\leq 
	\| f \|_2^2 e^{- \lambda_12  t}
	\end{align*}
	die Aussage.
\end{proof}




\section{Der W\"armeleitungskern unter Dirichletrandbedingung}

Sei $ G  $ ein Graph und $ S \leq G $ ein induzierter Teilgraph.
Wir betrachten den Laplaceoperator unter Dirichletrandbedingungen. Insbesondere erf\"ullen die Eigenfunktionen diese Randbedingung.
Durch analoges Vorgehen erhalten wir 
\begin{align}
i_\lambda^{D,S}(x,y)
=
\sum \limits_{\phi \in E_\lambda}
\phi(x) \phi(y).
\end{align}
Damit ist hierf\"ur automatisch die Dirichletrandbedingung erf\"ullt.
Falls $ S $ fixiert sein sollte, schreiben wir auch
$ i_\lambda^D := i_\lambda^{D,S} $.
Analog werden wir bei dem W\"armeleitungskern verfahren.

\begin{df}
	Sei $ t \in [0,\infty) $.
	Der W\"armeleitungsoperator $ H_t^D $ ist definiert durch
	\begin{align}
	H_t^D : L^2(S) \to L^2(S),
	u_0 \mapsto u(t,\cdot ),
	\end{align}
	wobei $ f \in  L^2(S) $ mit der Dirichletbedingung $ f(x) = 0 $ f\"ur $ x \in \delta S $ fortgesetzt wird.
	Hierbei ist $ u $ die L\"osung der W\"armeleitungsgleichung
	\eqref{eq:heat_equation} mit Dirichlet-Laplace.
\end{df}

 

\begin{thm}\label{th:repres_heat_kernel_dirichlet}
	F\"ur den W\"armeleitungsoperator gilt
	\begin{align}\label{eq:heat_kernel_def_dirichlet}
	u(t,x ) = (H_t^D u_0)(x) = 
	\sum \limits_{\lambda \in \sigma(\mathcal{L})} e^{- \lambda t } (I_\lambda u_0) (x) 
	=
	\langle u_0 ,h_t^D(x,\cdot) \rangle_S.
	%\sum \limits_{y \in V(G)} h_t(x,y) f(y) \mathrm{d}_y.
	\end{align}
	f\"ur die Anfangsbedingung $ u_0 \in L^2(S) $.
	Hierbei ist
	\begin{align}
	h_t^D(x,y) := \sum \limits_{\lambda \in \sigma(\mathcal{L})}
	e^{- \lambda t} i^D_\lambda(x,y) 
%	=
%	\begin{cases}
%	\sum \limits_{\lambda \in \sigma(\mathcal{L})}
%	e^{- \lambda t} i_\lambda(x,y) , & \  x \in S \wedge y \in S\\
%	\quad  \quad 0 , & \ \text{sonst}
%	\end{cases}
	\end{align}
	der W\"armeleitungskern auf dem Teilgraphen $ S $ mit Dirichletrandbedingung.
	Au\ss{}erdem bleibt die Randbedingung unabh\"angig von der Zeit erhalten.
\end{thm}

\begin{proof}
	Der Beweis verl\"auft analog zu Satz \ref{th:repres_heat_kernel_global}.
	Da die Eigenfunktionen die Dirichletrandbedingung erf\"ullen, gilt 
	\begin{align*}
	h_t^D(x,y) = 0
	\end{align*}
	f\"ur $ x \in \delta S$.  
	Damit folgt auch
	\begin{align}
	u(t,x) = \langle u_0, h_t^D(0,\cdot) \rangle_S  = 0
	\end{align}
	f\"ur $ x \in \delta S $ und $ t \in [0,\infty) $ beliebig.
\end{proof}

Wegen $ u(t,\cdot ) \in L^2(S) $ folgt mit dem Abschnitt 6.1  direkt
\begin{align}
\langle \mathcal{ L }u(t,\cdot) , u(t,\cdot ) \rangle_S
=
\mathcal{E}^S[u(t ,\cdot),u(t,\cdot)]. 
\end{align}
f\"ur den Dirichlet-Laplace.  


\begin{thm}\label{th:positivity_heat_kernel_dirichlet}
	Es gilt 
	\begin{align}
	h_t^D(x,y) \geq 0
	\end{align}
	f\"ur $ x,y \in S $.
\end{thm}

\begin{proof}
Der Beweis verl\"auft analog zu Satz \ref{th:pos_heat_kernel_global}. 
Der einzige Unterschied ist, dass $ v_\mathrm{min} $ in $ S $ bestimmt werden muss, da wir den Laplaceoperator auswerten wollen.

%	Sei $ f : S \cup \delta S  \to \mathbb{C}$ beliebig.
%	Dann existiert f\"ur ein beliebiges $ t \in [0, \infty) $ ein 
%	$ v_{\textrm{max}}  \in  S$ mit
%	\begin{align*}
%	(H_t f)(v_{max}) = F(v_{\mathrm{max}},t) = \max\limits_{v \in S } F(v,t).
%	\end{align*}
%	F\"ur den Laplace erhalten wir dann
%	\begin{align*}
%	\mathcal{L}f(v_{\mathrm{max}} ,t) =
%	\frac{1}{\mathrm{d}_{v_{\mathrm{max}}}}
%	\sum 
%	\limits_{u \sim v_{\mathrm{max}}}
%	(F(v_\mathrm{max},t) - F(u,t)) \geq 0
%	\end{align*}
%	Wegen \eqref{eq:fund_sol_2} folgt dann auch $ \partial_\mathrm{t} F(v_{\mathrm{max}},t) \leq 0 $.
%	Durch ein analoges Argument f\"ur das Minimum gilt $ \partial_\mathrm{t} F(v_\mathrm{min},t) \geq 0 $.
%	Damit wird insbesondere das Maximum mit der Zeit h\"ochstens kleiner und das Minimum gr\"o\ss{}er. 
%	Damit gilt f\"ur $ f \geq 0  $ auch $ (H_t f) \geq 0 $.
%	Mit 
%	\begin{align*}
%	f_z(y) =
%	\begin{cases}
%	\frac{1}{\mathrm{d}_z} , &\quad z=y\\
%	\ 0 ,						&\quad z \neq y
%	\end{cases} 
%	\end{align*}
%	erhalten wir
%	\begin{align*}
%	(H_t f )(x) = \sum \limits_{y \in S}
%	f(y) h_t(x,y) \mathrm{d}_y
%	= h_t(x,z) \geq 0
%	\end{align*}
%	für alle $ x,z \in S $.
%	Aufgrund der Symmetrie von $ h_t $ folgt die gew\"unschte Aussage.

\end{proof}

\begin{thm}
	Unter der Dirichletrandbedingung gilt
	\begin{align*}
	\| H_t^D f \|_2  \leq e^{- \lambda_1 t} \|f \|_2.
	\end{align*}
\end{thm}
\begin{proof}
	F\"ur die Dirichleteigenwerte gilt $ 0 < \lambda_1 \leq ....  \leq \lambda_n$.
	Da die Eigenfunktionen von $ \mathcal{L} $ eine ONB bilden gilt
	\begin{align*}
	\| H_t^D f \|_2^2
	\leq  
	\sum \limits_{i=1}^n 
	e^{-\lambda_i 2t}
	\left(\sum \limits_{x \in V(S)}  |(I_{\lambda_i} f) (x)|^2 \mathrm{d}_x\right)
	\leq 
	e^{-\lambda_1 2t} \| f \|_2^2.
	\end{align*}
\end{proof}

\begin{thm}
	Seien $ S \subset S^\prime $ induzierte Teilgraphen.
	Dann gilt
	\begin{align}
	h_t^{D,S}(x,y) \leq h_t^{D,S^\prime}(x,y)
	\end{align}
	für $ x,y \in S $.
\end{thm}

\section{Der W\"armeleitungskern unter Neumannrandbedingung}

%\textbf{\textcolor{red}{Kl\"aren wie Kern unter Neumannbed definiert ist.}}
Sei $ G  $ ein Graph und $ S \leq G $ ein induzierter Teilgraph.
Wir betrachten den Laplaceoperator unter Neumannrandbedingungen. Insbesondere erf\"ullen die Eigenfunktionen diese Randbedingung.
Durch analoges Vorgehen erhalten wir 
\begin{align}
i_\lambda^{N,S}(x,y)
=
\sum \limits_{\phi \in E_\lambda}
\phi(x) \phi(y).
\end{align}
Damit ist hierf\"ur automatisch die Neumannrandbedingung erf\"ullt.
Falls $ S $ fixiert sein sollte, schreiben wir auch
$ i_\lambda^N := i_\lambda^{N,S} $.
Analog werden wir bei dem W\"armeleitungskern verfahren.


\begin{df}
	Sei $ t \in [0,\infty) $.
	Der L\"osungsoperator $ H_t^N $ der W\"armeleitungsgleichung \eqref{eq:heat_equation} unter Neumannrandbedingung ist definiert durch
	\begin{align}
	H_t^N : L^2(S) \to L^2(S),
	u_0 \mapsto u(t,\cdot )
	\end{align}
	Insbesondere verwenden wir den Laplace unter Neumannrandbedingungen.
	Hierbei werden $ f \in L^2(S) $ mithilfe der Neumannrandbedingung
	\begin{align}
	\sum \limits_{\{x,y\} \in \partial S} f(x) - f(y) = 0 , \ x \in \delta S
	\end{align}
	eindeutig auf $ \delta S $ fortgesetzt.
\end{df}

Wir wollen zun\"achst die Fortsetzung genauer betrachten. 
Sei $ f : S \cup \delta S \to \R $ mit Neumannrandbedingung und $ x \in \delta S $.
Dann gilt
\begin{align*}
\sum \limits_{\{x,y\} \in \partial S} f(x) - f(y) = 0 , \ x \in \delta S
\
\Leftrightarrow
\
f(x) = \frac{1}{|x \sim S|} 
\sum 
\limits_{
	\substack{
		y \in S,\\
		y \sim x
		}
	}
f(y)
\end{align*}
mit $ x \sim S := \{ y \in S : x \sim y \} $.
Das bedeutet, dass die Randwerte bereits \"uber den Mittelwert der adjazenten Knoten bestimmt sind.\\


\begin{thm}\label{th:repres_heat_kernel_neumann}
	F\"ur den W\"armeleitungsoperator gilt
	\begin{align}\label{eq:heat_kernel_def_neumann}
	u(t,x ) = (H_t^N u_0)(x) = 
	\sum \limits_{\lambda \in \sigma(\mathcal{L})} e^{- \lambda t } (I_\lambda u_0) (x) 
	=
	\langle u_0, h_t^N(x,\cdot) \rangle_S.
	%\sum \limits_{y \in V(G)} h_t(x,y) f(y) \mathrm{d}_y.
	\end{align}
	f\"ur die Anfangsbedingung $ u_0 \in L^2(S) $.
	Insbesondere gilt $ u(0,x ) = 0 $.
	Hierbei ist
	\begin{align}
	h_t^N(x,y) :=
	\sum \limits_{\lambda \in \sigma(\mathcal{L})}
	e^{- \lambda t} i_\lambda^N(x,y)
%	
%	\begin{cases}
%	\sum \limits_{\lambda \in \sigma(\mathcal{L})}
%	e^{- \lambda t} i_\lambda(x,y) , & \  x \in S \wedge y \in S\\
%	\frac{1}{|x \sim S|}
%	\sum 
%	\limits_{
%		\substack{x \sim z
%			\\
%			z \in S
%			}
%		} h_t^N(z,y) , & \ x \in \delta S
%	\end{cases}
	\end{align}
	der W\"armeleitungskern auf dem Teilgraphen $ S $ mit Neumannrandbedingung.
	Au\ss{}erdem bleibt die Randbedingung unabh\"angig von der Zeit erhalten.
\end{thm}

\begin{proof}
	Der Beweis verl\"auft analog zu Satz \ref{th:repres_heat_kernel_global}.
	Da die Eigenfunktionen die Neumannrandbedingung erf\"ullen, gilt 
	\begin{align*}
	\sum \limits_{\{x,z \} \in \delta S} h_t^N(x,y) - h_t^N(z,y) = 0
	\end{align*}
	f\"ur $ x \in \delta S$.
	Damit erhalten wir auch
	\begin{align*}
	\sum \limits_{\{x,y \} \in \partial S} u(t,x) - u(t,y)
	=
	\sum \limits_{\{ x,y\} \in \partial S} \underbrace{\langle h_t(x,\cdot) - h_t(y,\cdot), u_0 \rangle_S}_{=0} = 0
	\end{align*}  
	f\"ur $ x \in \delta S $.
\end{proof}
 
Mit Abschnitt 6.2 erhalten wir 
\begin{align}
\langle \mathcal{ L }u(t,\cdot) , u(t,\cdot ) \rangle_S
=
\mathcal{E}^S[u(t ,\cdot),u(t,\cdot)]
\end{align}
f\"ur den Neumann-Laplace.


\begin{thm}
	Es gilt 
	\begin{align}
	h_t(x,y) \geq 0
	\end{align}
	f\"ur $ x,y \in S \cup \delta S $.
\end{thm}
\begin{proof}
	Wir bilden Maximum bzw. Minimum \"uber $ S \cup \delta S $. F\"ur die Auswertung des Laplace sollten sich diese im Inneren befinden.
	Angenommen $ x_\mathrm{max} \in \delta S $.
	Dann gilt 
	\begin{align*}
	\sum \limits_{
		\substack{x_\mathrm{max} \sim y \\ \ u \in S}
	}
	(u_0(x_\mathrm{max}) - u_0(y )) = 0.
	\end{align*}
	Wegen $ u_0(x_\mathrm{max}) \geq u_0(y) $ f\"ur alle $ y \in S $ muss zwangsl\"aufig $ u_0(x_\mathrm{max}) = f(y) $ f\"ur $ y \sim  x_\mathrm{max}$ gelten. 
	Damit befindet sich das Maximum immer auch im Inneren.
	Analog gilt dies f\"ur das Minimum.
	Der Rest ist analog zu den Satz \ref{th:pos_heat_kernel_global}.
\end{proof}


\begin{thm}
	Sei $ f \in L^2(S) $.
	Dann gilt
	\begin{align*}
	\| H_t f - \overline{f}  \|_2 \leq e^{-\lambda_1 t}. 
	\end{align*}
\end{thm}

\begin{proof}
Identisch zu \ref{th:L_2_estimate_global}.	
\end{proof}





\newpage
%%\section{Die Elementarl\"oesung der W\"armeleitungsgleichung}
%
%Wir betrachten von nun an einen beliebigen induzierten Subgraphen $ S $ von $ G $.
%
%\begin{thm}
%	Sei die Anfangsbedingung $ u_0 : S \cup \delta S \to \mathbb{C} $ gegeben.
%	Dann ist
%	\begin{align}\label{eq:fund_sol_1}
%	F(x,t)
%	:=
%	(H_t u_0)(x) =
%	\sum \limits_{y \in S \cup \delta S} h_t(x,y) u_0(y) \mathrm{d}_y
%	\end{align}
%	eine L\"osung der W\"armeleitungsgleichung mit $ F(x,0) = u_0(x) $.
%\end{thm}
%
%\begin{proof}
%	Wegen 
%	\begin{align*}
%	F(x,0 )=
%	e^{-0 \mathcal{L} } u_0(x) = u_0(x) 
%	\end{align*}
%	ist die Anfangsbedingung erf\"ullt und aufgrund von
%	\begin{align}\label{eq:fund_sol_2}
%	\frac{\partial}{\partial \mathrm{t}} F 
%	= 
%	\frac{\partial}{\partial \mathrm{t}} e^{-\mathcal{L} t} _0
%	=
%	- \mathcal{L} e^{-t \mathcal{L }} u_0 
%	= - \mathcal{L} F
%	\end{align}
%	ist $ F $ eine L\"osung der W\"armeleitungsgleichung.
%	\end{proof}
%
%Zur Erinnerung wiederholen wir die Dirichlet und-Neumannrandbedingung. Sei $ f : S \cup \delta S  \to \mathbb{R} $.
%Die Dirichletrandbedinung wird von $ f $ erf\"ullt, falls
%\begin{align}
%f(u) = 0
%\end{align}
%f\"ur alle $ u \in \delta S $ gilt. Analog wird die Neumannrandbedingung erf\"ullt, falls
%\begin{align}
%\sum 
%\limits_{\{u,v\} \in \partial S} f(u) - f(v) = 0
%\end{align} 
%f\"ur alle $ u \in \delta  $ gilt.
%\begin{thm}
%	Angenommen die Dirichlet -oder Neumannrandbedingung ist f\"ur $ f : S \cup \delta \mapsto \mathbb{R} $ erf\"ullt.
%	Dann ist diese auch f\"ur $ F $ unabh\"angig von der Zeit erf\"ullt.
%\end{thm}
%\begin{proof}
%	Wir nehmen an die Neumannrandbedingung ist erfüllt und zeigen die Aussage unter dieser Voraussetzung.
%	Wir wollen also
%	\begin{align*}
%	\sum \limits_{\{x,y\} \in \partial S} F(x,t) - F(y,t) = 0 
%	\end{align*}
%	f\"ur $ x \in \delta S $ zeigen.
%	Sei $ \lambda \in \sigma(\mathcal{L}) $ beliebig mit zugeh\"origer Eigenfunktion $ f_\lambda $, welche die Neumannrandbedingung erf\"ullt.
%	Dann gilt f\"ur $ x \in \delta S $
%	\begin{align*}
%	\sum \limits_{\{x,y\} \in \partial S} F(x,t) - F(y,t)
%	=
%	\sum \limits_{\{x,y\} \in \partial S} (H_t f) (x) - (H_t f)(y)
%	\overset{\ref{th:repres_heat_kernel}}{=}
%	e^{-\lambda t} \sum \limits_{\{x,y\} \in \partial S}  f (x) -  f(y) = 0.
%	\end{align*}
%	Dann gilt die Aussage insgesamt. \textcolor{red}{\textbf{Noch klären warum genau}}
%	F\"ur den Dirichletrand ist das Vorgehen analog.
%	
%\end{proof}
%
%Die Dirichletform f\"ur einen Teilgraph $ S \leq G $ ist durch
%\begin{align}
%\mathcal{E}^S [f,g]
%= 
%\sum \limits_{\{ u,v\} \in S^\ast} (f(u) - f(v)) \overline{(g(u) - g(v))}
%\end{align}
%f\"ur $ f,g \in \mathrm{Abb}(S \cup \delta S , \mathbb{C}) $ gegeben.
%\begin{thm}
%	Sei $ f : S \cup \delta S \to \mathbb{C} $ und die Neumann- oder Dirichletrandbedingung erf\"ullt.
%	Dann gilt
%	\begin{align}
%	\langle \mathcal{ L }F(\cdot,t) , F(\cdot ,t) \rangle_S
%	=
%	\mathcal{E}^S[F(\cdot ,t),F(\cdot,t)]. 
%	\end{align}
%	
%\end{thm}
%
%Im letzten Satz wurde gezeigt, dass sich die jeweiligen Randbedingungen von $ u_0 $ unabh\"angig  von $ t $ auf die L\"osung $ F $ \"ubertragen.
%Desweiteren wurde diese Identit\"at f\"ur beide Randtypen bereits in dem Vortrag von Arthur gezeigt.
%
%
%
%\section{Der W\"armeleitungskern unter Dirichletrand}
%
%In diesem Abschnitt sei $ f : S \cup \delta S \to \mathbb{C} $ mit Dirichletrand vorausgesetzt.
%
%
%\begin{thm}
%	Es gilt 
%	\begin{align}
%	h_t(x,y) \geq 0
%	\end{align}
%	f\"ur $ x,y \in S $.
%\end{thm}
%
%\begin{proof}
%	Sei $ f : S \cup \delta S  \to \mathbb{C}$ beliebig.
%	Dann existiert f\"ur ein beliebiges $ t \in [0, \infty) $ ein 
%	$ v_{\textrm{max}}  \in  S$ mit
%	\begin{align*}
%	(H_t f)(v_{max}) = F(v_{\mathrm{max}},t) = \max\limits_{v \in S } F(v,t).
%	\end{align*}
%	F\"ur den Laplace erhalten wir dann
%	\begin{align*}
%	\mathcal{L}f(v_{\mathrm{max}} ,t) =
%	\frac{1}{\mathrm{d}_{v_{\mathrm{max}}}}
%	\sum 
%	\limits_{u \sim v_{\mathrm{max}}}
%	(F(v_\mathrm{max},t) - F(u,t)) \geq 0
%	\end{align*}
%	Wegen \eqref{eq:fund_sol_2} folgt dann auch $ \partial_\mathrm{t} F(v_{\mathrm{max}},t) \leq 0 $.
%	Durch ein analoges Argument f\"ur das Minimum gilt $ \partial_\mathrm{t} F(v_\mathrm{min},t) \geq 0 $.
%	Damit wird insbesondere das Maximum mit der Zeit h\"ochstens kleiner und das Minimum gr\"o\ss{}er. 
%	Damit gilt f\"ur $ f \geq 0  $ auch $ (H_t f) \geq 0 $.
%	Mit 
%	\begin{align*}
%	f_z(y) =
%	\begin{cases}
%	\frac{1}{\mathrm{d}_z} , &\quad z=y\\
%	  \ 0 ,						&\quad z \neq y
%	\end{cases} 
%	\end{align*}
%	erhalten wir
%	\begin{align*}
%	(H_t f )(x) = \sum \limits_{y \in S}
%	f(y) h_t(x,y) \mathrm{d}_y
%	= h_t(x,z) \geq 0
%	\end{align*}
%	für alle $ x,z \in S $.
%	Aufgrund der Symmetrie von $ h_t $ folgt die gew\"unschte Aussage.
%\end{proof}
%
%\begin{thm}
%	Unter dem Dirichletrand gilt
%	\begin{align*}
%	\lim \limits_{t \to \infty}\| H_t f \|_2 = 0.
%	\end{align*}
%\end{thm}
%\begin{proof}
%	Für die Dirichleteigenwerte gilt $ 0 < \lambda_1 \leq ....  \leq \lambda_n$.
%	Da die Eigenfunktionen von $ \mathcal{L} $ eine ONB bilden gilt
%	\begin{align*}
%	\| H_t f \|_2
%	\leq  
%	\sum \limits_{i=1}^n 
%	e^{-\lambda_i 2t}
%	\left(\sum \limits_{x \in V(S)}  |(I_{\lambda_i} f) (x)|^2 \mathrm{d}_x\right)
%	=\sum \limits_{i=1}^n 
%	e^{-\lambda_i 2t}
%	\leq n e^{- \lambda_1 t}\overset{t \to \infty}{\rightarrow} 0. 
%	\end{align*}
%\end{proof}
%
%\section{Der W\"armeleitungskern unter Neumannrand}
%In diesem Abschnitt sei $ f : S \cup \delta S \to \mathbb{C} $ mit Neumannrand vorausgesetzt.
%Die hier bewiesenen Aussagen gelten ebenso auf dem Gesamtgraph $ G $.
%\begin{thm}
%	Es gilt 
%	\begin{align}
%	h_t(x,y) \geq 0
%	\end{align}
%	f\"ur $ x,y \in S \cup \delta S $.
%\end{thm}
%\begin{proof}
%	Der Beweis ist fast analog zu dem Dirichletfall.
%	Hier bilden wir jedoch Maximum bzw. Minimum über $ S \cup \delta S $. F\"ur die Auswertung des Laplace sollten sich diese im Inneren befinden.
%	Angenommen $ v_\mathrm{max} \in \delta S $.
%	Dann gilt 
%	\begin{align*}
%	\sum \limits_{
%		v_\mathrm{max} \sim u , \ u \in S
%		}
%		(f(v_\mathrm{max}) - f(u )) = 0.
%	\end{align*}
%	Wegen $ fG(v_\mathrm{max}) \geq f(u) $ f\"ur alle $ u \in S $ muss zwangsl\"aufig $ f(v_\mathrm{max}) = f(u) $ f\"ur $ u \sim  v_\mathrm{max}$ gelten. 
%	Damit befindet sich das Maximum immer auch im Inneren.
%	Analog gilt dies f\"ur das Minimum.
%\end{proof}
%
%\begin{thm}
%	Sei $ x \in S \cup \delta S $.
%	Dann gilt
%	\begin{align}
%	\sum \limits_{y \in S \cup \delta S}
%	h_t(x,y) d_y  = d_x.
%	\end{align}
%\end{thm}
%
%\begin{proof}
%	Wir betrachten die Abbildung
%	\begin{align*}
%	f : S \cup \delta S \to \mathbb{R}, \
%	x \mapsto \mathrm{d}_x.
%	\end{align*}
%	Damit gilt
%	\begin{align*}
%	(H_t f)(x) = \sum \limits_y h_t(x,y) \mathrm{d}_y
%	\end{align*}
%	und wegen $ 0 \in \sigma(\mathcal{L}) $ folgt
%	\begin{align*}
%	(H_t f)(x)
%	=
%	\sum \limits_{\lambda \in \sigma(\mathcal{L})}
%	e^{-\lambda t} I_\lambda f(x)
%	=
%	\mathrm{d}_x e^0 \textbf{1}(x) = \mathrm{d}_x
%	\end{align*}
%	und damit die Aussage.
%\end{proof}
% 
%
%\begin{thm}
%	Sei $ S  $ ein induzierter Teilgraph mit Neumannrand oder der gesamte Graph selbst.
%	Dann gilt
%	\begin{align*}
%	\lim 
%	\limits_{t \to \infty}
%	\| H_t f - I_0 f  \|_2 = 0. 
%	\end{align*}
%\end{thm}
%\begin{proof}
%	Wegen \ref{th:repres_heat_kernel} gilt
%	\begin{align*}
%	(H_t f)(x) - (I_0 f)(x) = 
%	\sum \limits_{\lambda \in \sigma(\mathcal{L}) \setminus \{0\}}
%	e^{-\lambda t} (I_\lambda f)(x)
%	=
%	\sum \limits_{i = 1}^{n-1} e^{-\lambda_i t} (I_{\lambda_i} f)(x)
%	\end{align*}
%	mit $ 0 < \lambda_1 < ... < \lambda_{n-1} $.
%	Damit erhalten wir durch
%	\begin{align*}
%	\| H_t f - I_0 f  \|_2 
%	&=
%	\sum \limits_{x \in V(S)} \left| 
%	\sum \limits_{i=1}^{n-1} e^{-\lambda_i t} I_{\lambda_i} f(x)
%	\right|^2 \mathrm{d}_x
%	\leq
%	\sum \limits_i e^{-\lambda_i t}  
%	\left(\sum \limits_x | I_{\lambda_i} f(x) |^2 \mathrm{d}_y\right)\\
%	&= 
%	\sum \limits_i e^{-\lambda_i t}
%	\leq 
%	(n-1 ) e^{-\lambda_1 t}
% 	\end{align*}
% 	die Aussage.
%
%\end{proof}
\end{document}
